\section*{Chapitre \ref{ch:b3:canonical-ensemble} --- L'ensemble canonique}

\begin{solution}{exo:b3:canonical-ensemble:1}
(a) $\mathcal P_2/\mathcal P_1 = (g_2/g_1)\eu^{-\beta\epsilon}$.
(b) $\beta\epsilon = 2.1/(8.62\times10^{-5} \times 2500) = 9.7$~:
fraction $3\eu^{-9.7} \approx \num{2e-4}$. (c) Une flamme porte
$\sim\num{e15}$ atomes de sodium par cm$^3$~: même $10^{-4}$ d'entre
eux, recyclés toutes les quelques nanosecondes, déversent $10^{19}$
photons jaunes par seconde --- de quoi éblouir. (d)
$3\eu^{-\beta\epsilon} = 0.1$~:
$T = \epsilon/(k_{\text{B}}\ln30) \approx \qty{7200}{K}$ --- une
photosphère stellaire, pas une flamme.
\end{solution}

\begin{solution}{exo:b3:canonical-ensemble:2}
(a) $z = 1 + \eu^{-\beta\epsilon} + \eu^{-2\beta\epsilon}$. (b)
$\langle E\rangle = \epsilon(\eu^{-\beta\epsilon} +
2\eu^{-2\beta\epsilon})/z$~: $\to 0$ à basse $T$, $\to \epsilon$
(le niveau moyen) à haute. (c) $\mathcal P_1 = 1/(\eu^{\beta
\epsilon} + 1 + \eu^{-\beta\epsilon})$ croît de façon monotone avec
$T$, n'approchant sa borne $1/3$ que pour $T \to \infty$. (d) Les
poids de Boltzmann ne font que décroître avec l'énergie à $T > 0$~;
une inversion de population exige les \omterm{ex:b3:microcanonical-ensemble:negative}{températures négatives} de
\cref{ex:b3:microcanonical-ensemble:negative}, inaccessibles à tout
réservoir.
\end{solution}

\begin{solution}{exo:b3:canonical-ensemble:3}
(a) Le facteur de Boltzmann sur $E_p = mgh$ à $T$ uniforme. (b) $h_0 =
k_{\text{B}}T/mg = \qty{8.4}{km}$. (c) $\eu^{-8848/8440} \approx
0.35$~: un tiers d'atmosphère --- d'où l'oxygène des sommiteux.
(d) L'atmosphère réelle n'est pas isotherme~: la température baisse
avec l'altitude, si bien que le profil s'écarte d'une exponentielle
unique.
\end{solution}

\begin{solution}{exo:b3:canonical-ensemble:4}
(a) Trois translations~: $\tfrac32 R = \qty{12.5}{J/(mol.K)}$, accord
exact. (b) On ajoute deux rotations~: $\tfrac52 R =
\qty{20.8}{J/(mol.K)}$, comme mesuré~: la vibration est gelée. (c)
$\tfrac72 R = \qty{29.1}{J/(mol.K)}$ prédits~; les $\approx26$
mesurés montrent une vibration à peine dégelée ($\theta_{\text{vib}}
\approx \qty{3400}{K}$). (d) $3R = \qty{24.9}{J/(mol.K)}$~:
Dulong--Petit, respectée à température ambiante par la plupart des
métaux et violée par le diamant --- l'histoire du gel de
\cref{exo:b3:canonical-ensemble:6}.
\end{solution}

\begin{solution}{exo:b3:canonical-ensemble:5}
(a) $z = \eu^{-\beta\hbar\omega/2}/(1 - \eu^{-\beta\hbar\omega})$.
(b) $\langle E\rangle = -\partial_\beta\ln z$ donne la forme annoncée
avec $\langle n\rangle = 1/(\eu^{\beta\hbar\omega} - 1)$. (c)
$C \to k_{\text{B}}$ à haute $T$~; $C \approx k_{\text{B}}(\beta
\hbar\omega)^2\eu^{-\beta\hbar\omega} \to 0$ à basse. (d) Avec
$\hbar\omega$ pour énergie d'un quantum de lumière, $\langle n\rangle$
est le nombre thermique de photons par mode~: la loi de Planck est à
un chapitre d'ici.
\end{solution}

\begin{solution}{exo:b3:canonical-ensemble:6}
(a) Numériquement, $C = \tfrac12 \times 3Nk_{\text{B}}$ vers
$T/\theta_{\text{E}} \approx 0.34$. (b) $\theta_{\text{E}}/T =
4.3$~: $C/3Nk_{\text{B}} = 4.3^2\eu^{-4.3}/(1 - \eu^{-4.3})^2
\approx 0.26$ --- le diamant à température ambiante n'a guère qu'un
quart de la capacité thermique classique~: l'«~anomalie~» est un gel
quantique en pleine lumière. (c) $\theta_{\text{E}} = \qty{90}{K}$
place le plomb en plein régime classique à \qty{300}{K}. (d)
L'hypothèse d'une fréquence unique~: les solides réels ont un spectre
de modes qui descend jusqu'au son de grande longueur d'onde, dont les
quanta bon marché donnent la loi en $T^3$ --- la réparation de Debye,
au \cref{ch:b3:photons-phonons}.
\end{solution}

\begin{solution}{exo:b3:canonical-ensemble:7}
(a) $f(v) \propto v^2\eu^{-mv^2/2k_{\text{B}}T}$~: $v_{\text{p}} =
\sqrt{2k_{\text{B}}T/m}$~: \qty{413}{m/s} pour N$_2$,
\qty{1540}{m/s} pour H$_2$. (b) $mv_{\text{esc}}^2/2k_{\text{B}}T
\approx 730$ pour N$_2$, $52$ pour H$_2$. (c) $\eu^{-730}$, c'est
jamais~; $\eu^{-52} \sim \num{e-23}$ par temps de séjour est lent ---
mais la haute atmosphère chaude ($\sim\qty{1000}{K}$) adoucit
l'exposant de l'hydrogène jusqu'à $\sim15$~: l'hydrogène s'échappe à
l'échelle géologique, l'azote reste. (d) La vitesse de libération et
la température de l'exosphère~: le puits à \qty{60}{km/s} de Jupiter
rend astronomique même l'exposant de l'hydrogène --- les géantes
gazeuses gardent ce que les petits mondes tièdes perdent.
\end{solution}

\begin{solution}{exo:b3:canonical-ensemble:8}
(a) $\partial_\beta^2\ln Z = \langle E^2\rangle - \langle
E\rangle^2$, et $C = \partial_T\langle E\rangle =
-k_{\text{B}}\beta^2\partial_\beta\langle E\rangle$. (b) Les deux
membres donnent $Nk_{\text{B}}(\beta\epsilon)^2\eu^{\beta\epsilon}/(
\eu^{\beta\epsilon} + 1)^2$. (c) $\langle E\rangle \propto N$,
$\Delta E \propto \sqrt N$. (d) L'intensité avec laquelle un système
répond au chauffage égale l'ampleur du tremblement spontané de son
énergie --- réponse et fluctuation sont deux lectures de la même
dérivée seconde, un motif (fluctuation--dissipation) qui revient
partout en physique.
\end{solution}

\begin{solution}{exo:b3:canonical-ensemble:9}
(a) $\ln2 = E_a\,\Delta T/k_{\text{B}}T^2$ avec $\Delta T =
\qty{10}{K}$, $T = \qty{298}{K}$~: $E_a =
0.693\,k_{\text{B}}T^2/10 \approx \qty{0.55}{eV}$. (b) De
$298$ à \qty{278}{K}~: facteur $\eu^{E_a(1/278 - 1/298)/k_{\text{B}}}
\approx 4.4$ de ralentissement --- d'où la conservation des aliments
au réfrigérateur. (c) $373 \to \qty{366}{K}$~: la vitesse chute de
$\approx1.4$~: l'œuf de onze minutes en demande un quart d'heure. (d)
Les réactions se gagnent sur la queue exponentielle des molécules
au-dessus de la barrière~: décaler un peu la température et la
population de cette queue bouge énormément --- la moyenne n'y est
presque pour rien.
\end{solution}

\begin{solution}{exo:b3:canonical-ensemble:10}
(a) $f_{\text{replié}} = 1/(1 + \eu^{-\beta(\Delta E -
T\Delta S)})$, l'entropie de l'état déplié étant absorbée dans son
\omterm{prop:b3:canonical-ensemble:machine}{énergie libre}. (b) En $T_{\text{m}} = \Delta E/\Delta S$ les deux
\omterm{prop:b3:canonical-ensemble:machine}{énergies libres} s'égalent~: moitié-moitié. (c) $T_{\text{m}} =
\num{4.8e-19}/\num{1.38e-21} = \qty{348}{K}$ (\qty{75}{\celsius})~;
largeur $\delta T \sim k_{\text{B}}T_{\text{m}}^2/\Delta E \approx
\qty{3.5}{K}$. (d) La coopérativité multiplie $\Delta E$ et
$\Delta S$ par le nombre de contacts qui lâchent ensemble, ce qui
laisse $T_{\text{m}}$ mais rétrécit la largeur $\propto 1/\Delta E$~: les
brins d'ADN se séparent sur deux kelvins, et c'est ce qui permet à une
machine PCR de cycler proprement entre «~fondu~» et «~apparié~».
\end{solution}

\begin{solution}{exo:b3:canonical-ensemble:11}
(a) $z = 2\cosh(\beta\mu B)$~; $M = N\mu\tanh(\beta\mu B) \approx
N\mu^2B/k_{\text{B}}T$ à petit argument~: le $1/T$ de Curie. (b)
$\mu_{\text{B}}B/k_{\text{B}}T$~: $\num{2.2e-3}$ à \qty{300}{K}~;
$0.67$ à \qty{1}{K} --- de l'indifférence à l'alignement marqué. (c)
$S$ ne dépend que de $\mu B/k_{\text{B}}T$~; baisser $B$ à entropie
fixée entraîne $T$ proportionnellement~: $\qty{1}{K} \to
\qty{0.1}{K}$. (d) Quand $k_{\text{B}}T$ atteint l'énergie
spin--spin, l'entropie n'est plus commandée par le champ~:
$\mu_0\mu_{\text{B}}^2/4\pi a^3 \approx \qty{7e-26}{J} \sim
\qty{5}{mK}$ --- le plancher classique (les moments nucléaires, mille
fois plus faibles, le repoussent au microkelvin).
\end{solution}

\begin{solution}{exo:b3:canonical-ensemble:12}
(a) $z_{\text{rot}} = \sum_J(2J+1)\eu^{-\beta BJ(J+1)}
\to \int(2J+1)\eu^{-\beta BJ(J+1)}\dd J = k_{\text{B}}T/B$~: alors
$\langle E\rangle = k_{\text{B}}T$ et $C_{\text{rot}} =
k_{\text{B}}$. (b) H$_2$~: $\theta_{\text{rot}} = \qty{87}{K}$ ---
la marche se situe dans la plage du laboratoire~; N$_2$~:
\qty{2.9}{K}, gelée seulement au voisinage de l'hélium liquide, si
bien que l'azote montre toujours $\tfrac52 R$. (c) Paliers
$\tfrac32 R$ (au-dessous de $\sim\qty{50}{K}$), $\tfrac52 R$ (de
$\sim200$ à $\sim\qty{1000}{K}$), puis montée vers $\tfrac72 R$ près de
$\theta_{\text{vib}} \approx
\qty{6000}{K}$. (d) Les $J$ pairs et
impairs appartiennent à des espèces de spin nucléaire différentes qui
s'interconvertissent lentement, si bien que le $C_V$ de l'hydrogène
froid dépend de son histoire ortho--para --- la statistique nucléaire
auditée au calorimètre.
\end{solution}

\begin{solution}{pb:b3:canonical-ensemble:1}
\textbf{1.} $V = \tfrac43\pi r^3 = \qty{4.0e-20}{m^3}$~: $m' = V
\Delta\rho = \qty{8.3e-18}{kg}$, poids $m'g = \qty{8.1e-17}{N}$.
\textbf{2.} $n(h) = n_0\,\eu^{-m'gh/k_{\text{B}}T}$~: la loi
barométrique, réduite à une lame de microscope.
\textbf{3.} $h_0 = k_{\text{B}}T/m'g = \num{4.04e-21}/\num{8.1e-17}
\approx \qty{50}{\micro m}$.
\textbf{4.} $h_0 \propto 1/r^3$~: une dispersion de $10\%$ sur le
rayon fait $30\%$ sur la hauteur caractéristique --- des grains
polydisperses brouillent l'échelle en bouillie. Perrin les a
fractionnés pendant des mois par centrifugations répétées.
\textbf{5.} Même formule, des masses distantes de $10^{10}$~:
l'«~atmosphère~» des grains est $10^{10}$ fois moins haute --- les
kilomètres se réduisent à quelques dizaines de micromètres, ce qui est
précisément ce qui la rend observable en entier au microscope.
\textbf{6.} Rapports successifs $0.55$, $0.55$, $0.57$~: constants à
l'erreur de comptage près --- exponentiel. $h_0 =
\qty{30}{\micro m}/\ln(100/55) \approx \qty{50}{\micro m}$.
\textbf{7.} $k_{\text{B}} = m'g\,h_0/T = \num{8.1e-17} \times
\num{5.0e-5}/293 \approx \qty{1.4e-23}{J/K}$.
\textbf{8.} $N_{\text{A}} = R/k_{\text{B}} \approx
\num{6.0e23}\,\unit{mol^{-1}}$.
\textbf{9.} À quelques pour cent près ici (avec des données
idéalisées)~; les séries réelles de Perrin se dispersaient de
$\pm15\%$ autour de la valeur moderne --- stupéfiant pour des grains
comptés à la main, et absolument décisif pour l'ordre de grandeur.
\textbf{10.} $m_{\text{H}} = \qty{e-3}{kg/mol}/N_{\text{A}} =
\qty{1.7e-27}{kg}$.
\textbf{11.} La chimie macroscopique fournit $R$~; la microscopie
optique et la pesée fournissent $m'$~; le comptage fournit
l'échelle~: trois mesures de paillasse triangulent la masse d'un atome
que personne ne peut voir.
\textbf{12.} Rien~: la démonstration n'a utilisé que «~système qui
échange de l'énergie avec un réservoir~» --- le facteur de Boltzmann
est aveugle à la taille, et c'est exactement ce que Perrin a vérifié.
\textbf{13.} $D = k_{\text{B}}T/6\pi\eta r =
\num{4.04e-21}/\num{4.0e-9} \approx \qty{e-12}{m^2/s}$.
\textbf{14.} $\sqrt{2Dt} = \sqrt{\num{1.2e-10}} \approx
\qty{11}{\micro m}$ par minute~: confortablement mesurable avec un
réticule d'oculaire et de la patience.
\textbf{15.} Deux phénomènes indépendants, deux formules
indépendantes, un seul nombre~: l'accord entre l'échelle statique et
le tremblement dynamique ne laissait aucune niche au hasard ---
l'hypothèse moléculaire prédisait les deux.
\textbf{16.} $\sqrt{k_{\text{B}}T/m} = \sqrt{\num{4.04e-21}/
\num{4.8e-17}} \approx \qty{9}{mm/s}$ par composante.
\textbf{17.} Le grain est heurté $10^{19}$ fois par seconde et oublie
sa vitesse en quelques nanomètres~: le vol balistique est
inobservable, et ce que l'œil voit est son intégrale --- la marche
aléatoire diffusive.
\textbf{18.} La traînée de Stokes (la viscosité des fluides du volume
de deuxième année) fournit le $6\pi\eta r$~; l'\omterm{thm:b3:canonical-ensemble:boltzmann}{ensemble canonique}
fournit le $k_{\text{B}}T$~: le $D$ d'Einstein est leur quotient, la
mécanique et la statistique dans une seule fraction.
\textbf{19.} Une fiction ne se compte pas~: dès lors que
$N_{\text{A}}$ est le rapport de deux nombres mesurés, avec des barres
d'erreur, les atomes deviennent des objets d'expérience --- Ostwald a
capitulé par écrit en 1909.
\textbf{20.} La diffusion qui bleuit le ciel, l'électrolyse jointe à
la charge mesurée de l'électron, l'hélium accumulé à partir du
radium~: quatre voies physiques sans rapport convergeant vers un même
$6\times10^{23}$ ont fait de ce nombre une propriété de la nature
plutôt que d'une théorie.
\textbf{21.} L'énergie gravitationnelle qui tire les grains vers le
bas, l'entropie configurationnelle qui les étale vers le haut,
échangées au taux $T$~: l'échelle \emph{est} le minimum de
$F = E - TS$.
\textbf{22.} $h_0 \propto 1/g$~: à $10^5g$,
$h_0 = k_{\text{B}}T/(m'\times10^5g) = \num{4.04e-21}/
\num{9.8e-17} \approx \qty{40}{\micro m}$ pour la protéine ---
l'équilibre de sédimentation, l'expérience de Perrin refaite chaque
jour dans les départements de biochimie.
\textbf{23.} De «~trop haute~» ($h_0 \gg$ profondeur du champ~:
aucun gradient visible) à «~trop mince~» ($h_0 \lesssim r$)~: les
rayons utilisables s'étendent grossièrement de $0.1$ à
$\qty{0.5}{\micro m}$ --- le choix de Perrin n'était pas de la chance
mais de la conception.
\textbf{24.} $k_{\text{B}}$ étant désormais exact par définition, la
même expérience étalonne les grains (leur taille ou leur masse
volumique) ou audite le montage~: la physique --- l'échelle de
Boltzmann --- est intacte~; seule a bougé la grandeur qui compte comme
inconnue.
\textbf{25.} La population d'un grain de \qty{0.2}{\micro m} est
divisée par deux tous les $\sim\qty{35}{\micro m}$~; lue à travers
$n_0\eu^{-m'gh/k_{\text{B}}
T}$, l'échelle a rendu $N_{\text{A}} \approx \num{6e23}$~: le nombre
d'Avogadro compté grain par grain --- et le débat atomique clos sur
une platine de microscope.
\end{solution}
